\chapter{Applications of Integration}

In the preceding chapters, we developed integration from Riemann sums to the
Fundamental Theorem, followed by techniques for finding antiderivatives---substitution,
integration by parts, partial fractions, and trigonometric methods---as well as
improper integrals for unbounded domains.

This chapter applies integration to geometric and physical quantities: areas
between curves, volumes of solids, arc lengths, surface areas, work, and
probability densities.

Underlying each application is a single idea:

\begin{center}
\emph{To integrate something is to add up infinitely many infinitely thin slices.}
\end{center}

Every problem here is solved the same way. You ask one question: \emph{what am I
accumulating?} You cut the thing you want into thin slices, you write down the
amount in one representative slice -- always of the form
$(\text{something})\,\dd x$ -- and then you let the integral sign add them all up.
The whole art is in finding the slice. The integral does the rest.

So this is not really a chapter of formulas. It is a chapter of \emph{one}
formula, $\int (\text{slice}) = \text{total}$, applied over and over until the
move becomes second nature. We begin where the integral itself began: with
area.

\chapterline

\section{The One Move: Slice, Measure, Sum}
\label{ch26:onemove}

Here is the three-step procedure we will use throughout the chapter.

\begin{focusbox}
\textbf{The slicing recipe.} To compute a total quantity $Q$ by integration:
\begin{enumerate}
  \item \textbf{Slice.} Cut the region (or solid, or curve, or process) into
  thin pieces indexed by a variable, say $x$, with thickness $\dd x$.
  \item \textbf{Measure.} Find the amount $\dd Q$ contributed by one slice. It
  will look like $\dd Q = (\text{some function of }x)\,\dd x$. This is the only
  creative step.
  \item \textbf{Sum.} Add the slices: $Q = \int_a^b \dd Q$, where $a$ and $b$
  are the smallest and largest values of the slicing variable.
\end{enumerate}
\end{focusbox}

Why does this work at all? Because that is the foundational meaning of the definite integral:
\[
  \int_a^b g(x)\,\dd x \;=\; \lim_{n\to\infty}\sum_{k=1}^n g(x_k^*)\,\Delta x,
\]
a limit of Riemann sums -- approximations formed by chopping an interval into $n$
subintervals of width $\Delta x = (b-a)/n$ and evaluating the height at a sample
point $x_k^*$. Each term $g(x_k^*)\,\Delta x$ is the contribution of
one thin strip of width $\Delta x$. When we write ``$\dd Q = g(x)\,\dd x$'' we are
naming a single term of that sum, and ``$\int \dd Q$'' is the instruction to take
the limit of the sum. So the slicing recipe is not a trick layered on top of
integration; it is integration, told as a story about pieces.

\begin{happeningbox}
When a physicist writes $\dd Q = g(x)\,\dd x$ and then ``integrates both sides,''
what is really happening is \emph{not} that $\dd x$ is a tiny number being
multiplied. The precise statement is: the total $Q$ is the limit of sums of
slice-contributions, and each slice-contribution is approximately $g(x_k)\,\Delta
x$, with the approximation getting perfect as $\Delta x \to 0$. The differential
notation $\dd Q = g(x)\,\dd x$ is just shorthand for ``the rate at which $Q$
accumulates per unit of $x$ is $g(x)$.'' The error you make in one slice by
treating it as a perfect rectangle is of \emph{higher order} -- it shrinks faster
than $\Delta x$ itself -- so it washes out completely in the limit. That higher
-order vanishing is the entire reason the recipe is exact rather than merely
approximate.
\end{happeningbox}

Keep that picture in mind. Now we make a slice, eight different ways.

\section{Area Between Curves}
\label{ch26:area}

The integral was born to measure the area under a single curve. The first
generalization is almost too easy -- the area \emph{between} two curves.

Suppose $f(x) \ge g(x)$ on $[a,b]$, so the graph of $f$ sits above the graph of
$g$. What is the area of the region trapped between them? Slice it vertically.

\begin{figure}[H]
\centering
\begin{tikzpicture}
\begin{axis}[
  width=10cm, height=6cm,
  xmin=-0.3, xmax=3.3, ymin=-0.5, ymax=5,
  axis lines=left, xlabel={$x$}, ylabel={$y$},
  xtick={0.5,2.6}, xticklabels={$a$,$b$},
  ytick=\empty, samples=100,
]
\addplot[name path=top, curveblue, thick, domain=0:3.2] {0.6*x^2+0.7};
\addplot[name path=bot, targetgreen, thick, domain=0:3.2] {0.9*x+0.2};
\addplot[areafill!55] fill between[of=top and bot, soft clip={domain=0.5:2.6}];
\addplot[curveblue] coordinates {(2.3,3.874)} node[above left,curveblue]{$f(x)$};
\addplot[targetgreen] coordinates {(2.6,2.54)} node[below right,targetgreen]{$g(x)$};
\fill[quotientorange!60] (axis cs:1.53,1.577) rectangle (axis cs:1.67,2.275);
\draw[black, thick] (axis cs:1.53,1.577) rectangle (axis cs:1.67,2.275);
\node[anchor=south east, font=\footnotesize] at (axis cs:1.35,2.7) {height $=f(x)-g(x)$};
\draw[->] (axis cs:1.35,2.65) -- (axis cs:1.52,2.0);
\draw[<->] (axis cs:1.53,0.15) -- (axis cs:1.67,0.15)
  node[midway,below,font=\tiny]{$\dd x$};
\draw[gray, dotted] (axis cs:1.53,0.15) -- (axis cs:1.53,1.577);
\draw[gray, dotted] (axis cs:1.67,0.15) -- (axis cs:1.67,1.703);
\end{axis}
\end{tikzpicture}
\caption{A vertical slice of the region between $f$ (top) and $g$ (bottom). One
slice is a thin rectangle of width $\dd x$ and height $f(x)-g(x)$.}
\label{ch26:fig-area}
\end{figure}

\textbf{Slice.} A vertical strip at position $x$, thickness $\dd x$.

\textbf{Measure.} The strip is (almost) a rectangle. Its width is $\dd x$ and its
height is the vertical gap between the curves, $f(x)-g(x)$. So
\[
  \dd A = \bigl(f(x)-g(x)\bigr)\,\dd x.
\]

\textbf{Sum.}
\[
  \boxed{\,A = \int_a^b \bigl(f(x)-g(x)\bigr)\,\dd x.\,}
\]

That is the entire theory. We did not memorize a new formula so much as
read it off the picture. The integrand is always (top minus bottom).

\begin{trapbox}
\textbf{Top minus bottom, not left limit minus right.} The single most common
error is forgetting which curve is on top, and getting a negative area. Area is
positive. If you are not sure which function is larger on the interval, test one
point: plug in a value of $x$ between $a$ and $b$ and see which output is bigger.
A second classic trap: if the curves \emph{cross} inside $[a,b]$, the top and
bottom switch. Then you must split the integral at the crossing points and use
$\abs{f-g}$ on each piece. Sketch first, always.
\end{trapbox}

\begin{example}[Area between a parabola and a line]
\label{ch26:ex-parab-line}
Find the area enclosed between $y = x^2$ and $y = x+2$.

First, where do they meet? Set $x^2 = x+2$, i.e. $x^2-x-2=0$, i.e.
$(x-2)(x+1)=0$. The curves cross at $x=-1$ and $x=2$. Between those points, which
is on top? Test $x=0$: the parabola gives $0$, the line gives $2$. The line is
on top. So
\[
  A = \int_{-1}^{2}\bigl[(x+2)-x^2\bigr]\,\dd x
    = \left[\frac{x^2}{2}+2x-\frac{x^3}{3}\right]_{-1}^{2}.
\]
At $x=2$: $2 + 4 - \tfrac{8}{3} = 6 - \tfrac{8}{3} = \tfrac{10}{3}$.
At $x=-1$: $\tfrac12 - 2 + \tfrac13 = \tfrac{3}{6}-\tfrac{12}{6}+\tfrac{2}{6}
= -\tfrac{7}{6}$.
Subtracting,
\[
  A = \frac{10}{3} - \left(-\frac{7}{6}\right) = \frac{20}{6}+\frac{7}{6}
    = \frac{27}{6} = \frac{9}{2}.
\]
The enclosed area is $\tfrac{9}{2}$.
\end{example}

\subsection{Slicing the other way}

Sometimes a region is easier to describe by horizontal strips -- when the
boundaries are naturally functions of $y$ rather than $x$. Then you slice
horizontally: a strip at height $y$ has thickness $\dd y$ and width
(right curve) minus (left curve) (Figure~\ref{ch26:fig-horiz-slice}), both expressed as functions of $y$. The area is
\[
  A = \int_c^d \bigl(x_{\text{right}}(y) - x_{\text{left}}(y)\bigr)\,\dd y.
\]
It is the same recipe with the roles of $x$ and $y$ swapped. Choose whichever
direction makes the slices simple.

\begin{figure}[H]
\centering
\begin{tikzpicture}[scale=1.05]
\draw[->] (-0.4,0)--(4.6,0) node[right]{$x$};
\draw[->] (0,-1.6)--(0,2.6) node[above]{$y$};
\fill[areafill!45] plot[domain=-1:2, samples=50] ({0.7*\x*\x}, {\x}) -- plot[domain=2:-1, samples=10] ({0.7*(\x+2)}, {\x}) -- cycle;
\draw[domain=-1.2:2.1, samples=60, curveblue, thick] plot ({0.7*\x*\x}, {\x});
\draw[domain=-1.3:2.1, samples=10, targetgreen, thick] plot ({0.7*(\x+2)}, {\x});
\fill[quotientorange!65] (0.252, 0.48) rectangle (1.82, 0.72);
\draw[black, thick] (0.252, 0.48) rectangle (1.82, 0.72);
\draw[<->] (-0.15, 0.48) -- (-0.15, 0.72) node[midway, left, font=\footnotesize] {$\dd y$};
\draw[<-] (1.82, 0.6) -- (2.25, 0.6) node[right, font=\footnotesize] {width $= x_{\text{right}} - x_{\text{left}}$};
\draw[dashed, gray] (-0.1, -1) node[left, black, font=\footnotesize] {$y=c$} -- (0.7, -1);
\draw[dashed, gray] (-0.1, 2) node[left, black, font=\footnotesize] {$y=d$} -- (2.8, 2);
\node[curveblue, left, font=\footnotesize] at (0.45, 1.35) {$x_{\text{left}}(y)$};
\node[targetgreen, right, font=\footnotesize] at (2.7, 1.7) {$x_{\text{right}}(y)$};
\end{tikzpicture}
\caption{A horizontal slice of a region bounded by curves $x = x_{\text{left}}(y)$ and $x = x_{\text{right}}(y)$. One slice is a thin horizontal rectangle of height $\dd y$ and width $x_{\text{right}}(y) - x_{\text{left}}(y)$.}
\label{ch26:fig-horiz-slice}
\end{figure}

\begin{example}[Same region, horizontal slices]
\label{ch26:ex-horiz}
Find the area enclosed by $x = y^2$ and $x = y+2$.

Solving $y^2 = y+2$ gives $y^2-y-2 = (y-2)(y+1) = 0$, so $y$ runs from $-1$ to
$2$. For these $y$, the line $x=y+2$ is to the right of the parabola $x=y^2$
(check $y=0$: line gives $2$, parabola gives $0$). So
\[
  A = \int_{-1}^{2}\bigl[(y+2)-y^2\bigr]\,\dd y = \frac{9}{2},
\]
identical to before -- it is literally the same algebra. The lesson: pick the
slicing direction that makes ``right minus left'' (or ``top minus bottom'') a
single clean expression, so you never have to split the integral.
\end{example}

\begin{walkbox}
Picture a region shaped like a sideways parabola, opening to the right, capped by
a vertical line. If you slice it vertically, a single vertical strip can hit the
parabola \emph{twice} (top and bottom branches), forcing you to split the
integral. Slice it horizontally and each strip hits one left edge and one right
edge cleanly. Which way would you cut, and why does the geometry of the region --
not any rule -- decide for you?
\end{walkbox}

\section{Volumes by Slicing}
\label{ch26:slicing}

Now lift the idea into three dimensions. The recipe does not change at all: cut
the solid into thin slices, find the volume of one slice, integrate.

The cleanest case is slicing a solid by parallel planes. Imagine a loaf of bread.
Cut it into thin slices perpendicular to the $x$-axis. A slice at position $x$ is
a flat slab of thickness $\dd x$ whose flat face has some area $A(x)$. Its volume
is (face area)$\times$(thickness):
\[
  \dd V = A(x)\,\dd x, \qquad\text{so}\qquad
  \boxed{\,V = \int_a^b A(x)\,\dd x.\,}
\]
This is the general principle of volume by slicing, and \emph{everything} that
follows -- disks, washers, even shells -- is a special case of it. You only ever
need to know the cross-sectional area $A(x)$ as a function of where you cut.

\begin{figure}[H]
\centering
\begin{tikzpicture}[scale=1.0]
\draw[->] (-0.4,0)--(7.2,0) node[right]{$x$};
\foreach \x/\h in {0.6/0.55, 1.4/0.95, 2.2/1.25, 3.0/1.4, 3.8/1.35, 4.6/1.1, 5.4/0.8, 6.2/0.45}{
  \draw[curveblue, fill=areafill!35] (\x,-\h) rectangle (\x+0.18,\h);
}
\draw[black, very thick, fill=areafill!70] (3.0,-1.4) rectangle (3.18,1.4);
\node[above, font=\footnotesize] at (3.09,1.7) {face area $A(x)$};
\draw[->, thin] (3.09,1.65) -- (3.09,1.42);
\draw[<->] (3.0,-1.75)--(3.18,-1.75);
\node[below, font=\footnotesize] at (3.09,-1.8) {$\dd x$};
\end{tikzpicture}
\caption{A solid cut into slabs perpendicular to the $x$-axis. One slab has
volume $A(x)\,\dd x$, where $A(x)$ is the area of its face.}
\label{ch26:fig-slicing}
\end{figure}

\begin{example}[A pyramid, derived from scratch]
\label{ch26:ex-pyramid}
A pyramid has a square base of side $s$ and height $h$, apex directly above the
center. Find its volume.

Put the apex at the origin and let $x$ measure distance \emph{down} from the apex,
$0 \le x \le h$. At depth $x$ the cross-section is a square. By similar triangles,
the square's side grows linearly from $0$ at the apex to $s$ at the base:
side $= s\,(x/h)$. So the cross-sectional area is
\[
  A(x) = \left(\frac{s x}{h}\right)^2 = \frac{s^2}{h^2}\,x^2.
\]
Therefore
\[
  V = \int_0^h \frac{s^2}{h^2}\,x^2\,\dd x
    = \frac{s^2}{h^2}\cdot\frac{x^3}{3}\Big|_0^h
    = \frac{s^2}{h^2}\cdot\frac{h^3}{3}
    = \frac{1}{3}s^2 h.
\]
This recovers the familiar formula: the volume of a pyramid is one-third base times height.
The factor $\tfrac13$ from elementary geometry arises directly from $\int_0^h x^2\,\dd x = \tfrac{h^3}{3}$;
because the cross-sectional area scales quadratically with depth, integrating $x^2$ produces the factor of $\tfrac13$.
\end{example}

\begin{historybox}
The factor of $\tfrac13$ in the volume of a pyramid has a long history. In the 5th
century BCE, Democritus conjectured that a pyramid has one-third the volume of a
prism with the same base and height. Eudoxus (c.~370 BCE) and Archimedes (c.~250 BCE)
proved this using the \emph{method of exhaustion}: trapping the solid between inner
and outer approximations of prisms and showing the difference vanishes in the limit.
Archimedes also determined the volume of a sphere by balancing slices against a cone
and cylinder.

In 1615, Johannes Kepler estimated the capacity of wine barrels by slicing them into
thin circular disks. In 1635, Bonaventura Cavalieri generalized this into
\emph{Cavalieri's principle}: two solids have the same volume if all horizontal
cross-sections at equal heights have equal areas. Cavalieri viewed a solid as a stack
of two-dimensional planes. Critics questioned how zero-thickness planes could form a
nonzero volume, but the method gave consistent results.

Calculus resolved this by replacing indivisible planes with differential slices:
$\dd V = A(x)\,\dd x$ represents a rate of volume accumulation, integrated via an
antiderivative.
\end{historybox}

\subsection{Solids of revolution: the disk method}

Now a very common special case: spin a region around an axis to sweep out a solid.
If you revolve the region under $y=f(x)$ (above the $x$-axis, from $a$ to $b$)
about the $x$-axis, each vertical strip sweeps out a thin \emph{disk} -- a coin --
of radius $f(x)$ and thickness $\dd x$. A disk is just a cylinder, so its area is
$\pi(\text{radius})^2$:
\[
  A(x) = \pi\,[f(x)]^2, \qquad
  \boxed{\,V = \pi\int_a^b [f(x)]^2\,\dd x.\,}
\]
This is not a new method -- it is volume-by-slicing with a circular cross-section.

\begin{figure}[H]
\centering
\begin{tikzpicture}
\begin{axis}[
  width=10cm, height=5cm,
  xmin=-0.3, xmax=3.4, ymin=-2.2, ymax=2.2,
  axis lines=middle, xlabel={$x$}, ylabel={$y$},
  xtick=\empty, ytick=\empty, samples=80,
]
\addplot[curveblue, thick, domain=0:3.2] {0.5*x+0.4};
\addplot[curveblue, thick, domain=0:3.2] {-(0.5*x+0.4)};
\draw[black, fill=areafill!50] (axis cs:1.6,0) ellipse [x radius=3pt, y radius=34pt];
\draw[tangentred, very thick] (axis cs:1.6,0) -- (axis cs:1.6,1.2) node[midway, left, font=\footnotesize] {$f(x)$};
\fill[black] (axis cs:1.6,0) circle (1.5pt);
\node[below, font=\footnotesize] at (axis cs:1.6,-1.4) {disk of thickness $\dd x$};
\end{axis}
\end{tikzpicture}
\caption{Revolving $y=f(x)$ about the $x$-axis. A vertical strip sweeps out a disk
of radius $f(x)$; its volume is $\pi[f(x)]^2\,\dd x$.}
\label{ch26:fig-disk}
\end{figure}

\begin{example}[Volume of a sphere -- Archimedes' tombstone]
\label{ch26:ex-sphere}
Revolve the upper half of the circle $y = \sqrt{r^2 - x^2}$, from $x=-r$ to
$x=r$, about the $x$-axis. You sweep out a solid ball of radius $r$. The disk at
$x$ has radius $\sqrt{r^2-x^2}$, so
\[
  V = \pi\int_{-r}^{r}\bigl(\sqrt{r^2-x^2}\bigr)^2\,\dd x
    = \pi\int_{-r}^{r}(r^2 - x^2)\,\dd x.
\]
The integrand is even, so integrate from $0$ to $r$ and double:
\[
  V = 2\pi\left[r^2 x - \frac{x^3}{3}\right]_0^{r}
    = 2\pi\left(r^3 - \frac{r^3}{3}\right)
    = 2\pi\cdot\frac{2r^3}{3} = \frac{4}{3}\pi r^3.
\]
This derives the classical sphere volume $\tfrac43\pi r^3$ directly.
\end{example}

\subsection{Hollow solids: the washer method}

What if the region you revolve does not touch the axis -- if there is a gap? Then
each slice is not a solid disk but a \emph{washer}: a disk with a hole punched out.
Revolve the region between $y=f(x)$ (outer) and $y=g(x)$ (inner), with
$f \ge g \ge 0$, about the $x$-axis. The slice at $x$ is an \emph{annulus} --
a flat ring, like a metal washer or a DVD, bounded by two concentric circles --
with outer radius $f(x)$ and inner radius $g(x)$ (Figure~\ref{ch26:fig-washer}). Its area is
the big circle minus the little one:
\[
  A(x) = \pi[f(x)]^2 - \pi[g(x)]^2,\qquad
  \boxed{\,V = \pi\int_a^b\Bigl([f(x)]^2 - [g(x)]^2\Bigr)\,\dd x.\,}
\]

\begin{figure}[H]
\centering
\begin{tikzpicture}
\fill[areafill!55] (0,0) circle (1.5cm);
\fill[white] (0,0) circle (0.7cm);
\draw[curveblue, very thick] (0,0) circle (1.5cm);
\draw[targetgreen, very thick] (0,0) circle (0.7cm);
\draw[->, thick] (0,0)--(1.5cm,0); \node[above] at (0.75cm,0) {\footnotesize $f(x)$};
\draw[->, thick] (0,0)--(-0.7cm,0); \node[above] at (-0.35cm,0) {\footnotesize $g(x)$};
\node[curveblue] at (1.15cm,1.15cm) {\footnotesize outer};
\node[targetgreen, below] at (0,-0.78cm) {\footnotesize hole};
\end{tikzpicture}
\caption{A washer cross-section: a disk of outer radius $f(x)$ with a concentric
hole of inner radius $g(x)$ punched out. Its area is $\pi[f(x)]^2 - \pi[g(x)]^2$.}
\label{ch26:fig-washer}
\end{figure}

\begin{trapbox}
\textbf{Square first, then subtract -- never the other way.} The washer area is
$\pi(R^2 - r^2)$, \emph{not} $\pi(R-r)^2$. The annulus is the difference of two
\emph{areas}, so you square each radius and then subtract. Writing
$\pi\int (f-g)^2$ is a wrong answer that looks plausible; it computes a solid of a
completely different shape. Always: outer radius squared minus inner radius
squared.
\end{trapbox}

\begin{example}[A washer volume]
\label{ch26:ex-washer}
Revolve the region between $y=x$ and $y=x^2$ (for $0\le x\le 1$) about the
$x$-axis. On this interval $x \ge x^2$, so the outer radius is $f(x)=x$ and the
inner radius is $g(x)=x^2$. Then
\[
  V = \pi\int_0^1\bigl(x^2 - x^4\bigr)\,\dd x
    = \pi\left[\frac{x^3}{3} - \frac{x^5}{5}\right]_0^1
    = \pi\left(\frac13 - \frac15\right)
    = \pi\cdot\frac{2}{15} = \frac{2\pi}{15}.
\]
\end{example}

\subsection{The shell method: slicing parallel to the axis}

There is a second way to slice a solid of revolution, and it is worth
seeing because it pushes the slicing idea a level deeper. Instead of
cutting \emph{across} the axis into disks, cut \emph{around} the axis into nested
cylindrical shells -- think of the rings of a tree, or a set of nesting tin cans.

Revolve the region under $y=f(x)$, from $x=a$ to $x=b$, about the \emph{$y$-axis}.
Take a thin vertical strip at position $x$, width $\dd x$, height $f(x)$. When this
strip revolves around the $y$-axis it sweeps out a thin cylindrical shell of radius
$x$ and height $f(x)$.

Now, what is the volume of a thin shell? \emph{Unroll
it.} A thin cylindrical shell of radius $x$, height $f(x)$, and thickness $\dd x$,
slit open and flattened, is a thin rectangular slab. Its dimensions are:
\begin{itemize}
  \item length = circumference of the shell = $2\pi x$,
  \item height = $f(x)$,
  \item thickness = $\dd x$.
\end{itemize}
So its volume is
\[
  \dd V = \underbrace{2\pi x}_{\text{circumference}}\cdot
          \underbrace{f(x)}_{\text{height}}\cdot
          \underbrace{\dd x}_{\text{thickness}},
\]
and summing the shells,
\[
  \boxed{\,V = \int_a^b 2\pi x\,f(x)\,\dd x.\,}
\]

\begin{figure}[H]
\centering
\begin{tikzpicture}
\begin{scope}
\draw[->] (0,0)--(4,0) node[right]{$x$};
\draw[->] (0,0)--(0,3) node[above]{$y$};
\draw[curveblue, thick, domain=0:3.4, smooth, samples=40]
  plot (\x, {2.6*exp(-0.35*\x)});
\draw[black, very thick, fill=areafill!50] (2.0,0) rectangle (2.18,1.29);
\node[below] at (2.1,0) {\footnotesize $x$};
\draw[<->] (2.0,-0.3)--(2.18,-0.3); \node[below] at (2.1,-0.32){\tiny $\dd x$};
\node at (1.3,2.6) {\footnotesize $y=f(x)$};
\end{scope}
\begin{scope}[xshift=5.6cm]
\node at (1.6,3.0) {\footnotesize unrolled shell};
\draw[black, fill=areafill!50] (0,0) rectangle (3.2,1.29);
\draw[<->] (0,-0.3)--(3.2,-0.3);
\node[below] at (1.6,-0.32) {\footnotesize $2\pi x$};
\draw[<->] (-0.3,0)--(-0.3,1.29);
\node[left] at (-0.3,0.65) {\footnotesize $f(x)$};
\end{scope}
\end{tikzpicture}
\caption{The shell method. A vertical strip at $x$ revolves into a cylindrical
shell; unrolled, the shell is a slab of dimensions $2\pi x$ by $f(x)$ by $\dd x$.}
\label{ch26:fig-shell}
\end{figure}

\begin{whybox}
\textbf{Why ``unrolling'' is legitimate.} A real cylindrical shell is not a flat
slab -- its outer face ($2\pi(x+\dd x)$ around) is slightly longer than its inner
face ($2\pi x$ around). So why may we use the single circumference $2\pi x$? Write
the exact shell volume as the difference of two cylinders:
\[
  \dd V = \pi(x+\dd x)^2 f(x) - \pi x^2 f(x)
        = \pi f(x)\bigl(2x\,\dd x + (\dd x)^2\bigr)
        = 2\pi x\,f(x)\,\dd x + \pi f(x)(\dd x)^2.
\]
The leading term is exactly $2\pi x\,f(x)\,\dd x$. The correction
$\pi f(x)(\dd x)^2$ involves $(\dd x)^2$ -- it is \emph{higher order}, shrinking
faster than $\dd x$ -- so when we sum and take the limit it contributes nothing.
This is the same higher-order-vanishing that justified treating an area strip as a
perfect rectangle. Unrolling is legitimate because the discrepancy between the inner
and outer circumferences is a second-order effect.
\end{whybox}

\begin{example}[Same solid, two methods]
\label{ch26:ex-shell-check}
Revolve the region under $y = x^2$, from $x=0$ to $x=2$, about the $y$-axis.

\textbf{By shells} (slicing parallel to the $y$-axis):
\[
  V = \int_0^2 2\pi x\cdot x^2\,\dd x = 2\pi\int_0^2 x^3\,\dd x
    = 2\pi\cdot\frac{x^4}{4}\Big|_0^2 = 2\pi\cdot 4 = 8\pi.
\]

\textbf{By washers}, to check, we slice perpendicular to the $y$-axis. Here $y$
runs from $0$ to $4$; at height $y$ the solid is the region between the curve
$x=\sqrt{y}$ and the line $x=2$, revolved, giving outer radius $2$ and inner
radius $\sqrt{y}$:
\[
  V = \pi\int_0^4\bigl(2^2 - (\sqrt y)^2\bigr)\,\dd y
    = \pi\int_0^4 (4 - y)\,\dd y
    = \pi\left[4y - \frac{y^2}{2}\right]_0^4
    = \pi(16-8) = 8\pi.
\]
The two methods agree, as they must. Notice how shells here avoided having to
solve for $x$ in terms of $y$. \emph{That} is when you reach for shells: choose
the method whose slices are easy to describe.
\end{example}

\begin{intuitionbox}
A pocket rule for choosing. Ask: in which direction are my slices parallel to the
axis of rotation? Slices \emph{perpendicular} to the axis are disks/washers; slices
\emph{parallel} to the axis are shells. Then ask which description of the boundary
is simpler. Revolving around the $y$-axis a region given as $y=f(x)$? Vertical
strips are parallel to the axis, so shells let you keep the natural variable $x$ --
no algebraic gymnastics. The methods never disagree; one is just less work.
\end{intuitionbox}

\subsection{The Pappus--Guldin Theorems: Revolution via Centroids}

\begin{connectionbox}
Before calculus gave us disks and shells, the ancient Greek geometer \emph{Pappus of Alexandria}
(c.~320 CE, in Book~VII of his \emph{Collection}) discovered an astonishing shortcut for solids
of revolution, later rediscovered and popularized by the Swiss mathematician \emph{Paul Guldin} in 1640.

The \textbf{Pappus--Guldin Centroid Theorems} state that when a plane region of area $A$ is revolved
about an external axis in its plane (an axis lying in the same plane that touches at most the boundary
of the region, never cutting through its interior), the volume swept out is simply the area times
the distance traveled by its geometric center of mass -- its \emph{centroid} $\bar{y}$, the point
where the flat shape would balance on the tip of a pin:
\[
  V = 2\pi \bar{y} A.
\]
Similarly, the surface area generated by revolving a curve of arc length $L$ about an external axis is
the length times the circumference traced by the curve's centroid:
\[
  S = 2\pi \bar{y}_{\text{curve}} L.
\]
Why does this work? In the shell method, the volume integral is $\int 2\pi x f(x)\,\dd x = 2\pi \int x f(x)\,\dd x$.
By definition, the horizontal coordinate of the centroid of the region is $\bar{x} = \frac{1}{A}\int x f(x)\,\dd x$,
so the integral is identically $2\pi \bar{x} A$. Slicing by shells reveals that Guldin's theorem
directly expresses the shell integral in terms of the centroid of the planar region.
\end{connectionbox}

\section{Arc Length: Accumulating Distance}
\label{ch26:arclength}

So far we have accumulated area and volume. Now accumulate \emph{length}. How long
is a curve $y=f(x)$ from $x=a$ to $x=b$?

\textbf{Slice.} Cut the curve into tiny pieces. Over a sliver of horizontal width
$\dd x$, the curve rises by $\dd y$. The piece of curve is so short that it is
essentially straight -- a tiny line segment (smooth means locally straight, the
refrain of this whole book).

\textbf{Measure.} That little segment is the hypotenuse of a right triangle with
legs $\dd x$ and $\dd y$. Pythagoras gives its length:
\[
  \dd s = \sqrt{(\dd x)^2 + (\dd y)^2}.
\]
Factor out $\dd x$ to express everything in terms of the slope:
\[
  \dd s = \sqrt{1 + \left(\frac{\dd y}{\dd x}\right)^2}\;\dd x
        = \sqrt{1 + [f'(x)]^2}\,\dd x.
\]

\textbf{Sum.}
\[
  \boxed{\,L = \int_a^b \sqrt{1 + [f'(x)]^2}\,\dd x.\,}
\]

\begin{figure}[H]
\centering
\begin{tikzpicture}
\begin{axis}[
  width=9cm, height=5cm,
  xmin=0, xmax=3.4, ymin=0, ymax=3,
  axis lines=left, xlabel={$x$}, ylabel={$y$},
  xtick=\empty, ytick=\empty, samples=80,
]
\addplot[curveblue, very thick, domain=0.2:3.2] {0.25*x^2+0.4};
\coordinate (P) at (axis cs:1.6,1.04);
\coordinate (Q) at (axis cs:2.2,1.61);
\draw[black, very thick] (P)--(Q);
\draw[gray, dashed] (P)--(axis cs:2.2,1.04)--(Q);
\node at (axis cs:1.9,0.92) {\footnotesize $\dd x$};
\node at (axis cs:2.34,1.32) {\footnotesize $\dd y$};
\node at (axis cs:1.74,1.42) {\footnotesize $\dd s$};
\end{axis}
\end{tikzpicture}
\caption{A short piece of curve is the hypotenuse of a triangle with legs $\dd x$
and $\dd y$. Its length is $\dd s = \sqrt{(\dd x)^2+(\dd y)^2}$.}
\label{ch26:fig-arclength}
\end{figure}

\begin{happeningbox}
The symbol $\dd s = \sqrt{(\dd x)^2 + (\dd y)^2}$ is the infinitesimal version of
the distance formula. What is really happening: we replace the curve by an
inscribed polygon -- a chain of straight chords -- and add up the chord lengths.
Each chord has length $\sqrt{(\Delta x)^2 + (\Delta y)^2}$. As the chords get
shorter the polygon hugs the curve, and the sum of chord lengths approaches the
arc length. The formula is just that limit written compactly. Arc length is the
total ``ground covered,'' integrating the local stretch factor
$\sqrt{1+(f')^2}$ that tells you how much real distance you travel per unit of
horizontal travel.
\end{happeningbox}

\begin{trapbox}
\textbf{The square root usually fights back.} Arc-length integrands are
notoriously ugly: $\sqrt{1+[f'(x)]^2}$ rarely has an elementary antiderivative.
Even for $y=x^2$ you get $\int\sqrt{1+4x^2}\,\dd x$, which needs a trig or
hyperbolic substitution. Authors of textbook problems rig $f$ so the radical
collapses to a perfect square. In real life you reach for numerical integration.
Do not be alarmed when an arc-length integral looks intractable -- that is the
normal state of affairs, not a sign you set it up wrong.
\end{trapbox}

\begin{example}[A length that actually works out]
\label{ch26:ex-arclen}
Find the arc length of $y = \tfrac{2}{3}x^{3/2}$ from $x=0$ to $x=3$. Then
$f'(x) = x^{1/2}$, so $[f'(x)]^2 = x$, and
\[
  L = \int_0^3 \sqrt{1+x}\,\dd x
    = \frac{2}{3}(1+x)^{3/2}\Big|_0^3
    = \frac{2}{3}\bigl(4^{3/2} - 1\bigr)
    = \frac{2}{3}(8-1) = \frac{14}{3}.
\]
The radical collapsed because $1+(f')^2 = 1+x$ is itself a clean power. That is the
only reason the problem is doable by hand -- and it is exactly the kind of luck
textbook authors arrange.
\end{example}

\begin{recoverybox}
\textbf{Pause and think: does the formula work for a straight line?}
Consider a line $y = mx + b$ on $[a,b]$. The derivative is constant: $f'(x) = m$.
Plugging this into the arc length formula gives
\[
  L = \int_a^b \sqrt{1 + m^2}\,\dd x = \sqrt{1 + m^2}\,(b - a).
\]
Now look at the right triangle under the segment: its horizontal run is $\Delta x = b - a$
and its vertical rise is $\Delta y = m(b - a)$. The Pythagorean theorem gives the hypotenuse
as $\sqrt{(\Delta x)^2 + (\Delta y)^2} = \sqrt{(b-a)^2 + m^2(b-a)^2} = \sqrt{1+m^2}\,(b-a)$.
The calculus formula gives the exact Pythagorean length! The integral does not replace
geometry; it applies ordinary Pythagoras slice-by-slice where the slope $f'(x)$ is constantly bending.
\end{recoverybox}

\begin{connectionbox}
The little relation $\dd s = \sqrt{(\dd x)^2 + (\dd y)^2}$ is the seed of an
enormous subject. In physics, dividing by $\dd t$ gives speed
$\frac{\dd s}{\dd t} = \sqrt{\dot x^2 + \dot y^2}$ -- the magnitude of the velocity
vector. In differential geometry, the rule for measuring $\dd s$ at each point is
called the \emph{metric}, and replacing the flat Pythagorean $\dd s^2 = \dd x^2 +
\dd y^2$ with a position-dependent rule is precisely how Einstein's general
relativity encodes curved spacetime. Arc length is your first encounter with a
metric.
\end{connectionbox}

\section{Surface Area of Revolution}
\label{ch26:surface}

Combine the last two ideas. Revolve the curve $y=f(x)$ about the $x$-axis and ask
for the \emph{area of the surface} it sweeps out -- the skin of the solid, not its
volume.

\textbf{Slice.} A tiny arc of the curve, of length $\dd s$, at horizontal position
$x$ and height $f(x)$.

\textbf{Measure.} When that little arc revolves around the $x$-axis it traces a
thin \emph{band} -- a ribbon (Figure~\ref{ch26:fig-surface}). The ribbon is a tiny
strip wrapped in a circle of radius $f(x)$ (the distance from the curve to the
axis). Its circumference is
$2\pi f(x)$ and its width is the arc length $\dd s$. So its area is
\[
  \dd S = 2\pi\,f(x)\,\dd s = 2\pi f(x)\sqrt{1+[f'(x)]^2}\,\dd x.
\]

\textbf{Sum.}
\[
  \boxed{\,S = \int_a^b 2\pi\,f(x)\sqrt{1+[f'(x)]^2}\,\dd x.\,}
\]

\begin{figure}[H]
\centering
\begin{tikzpicture}
\begin{axis}[
  width=9cm, height=5.2cm,
  xmin=-0.3, xmax=3.4, ymin=-2.3, ymax=1.9,
  axis lines=middle, xlabel={$x$}, ylabel={$y$},
  xtick=\empty, ytick=\empty, samples=80,
]
\addplot[curveblue, thick, domain=0:3.2] {0.45*x+0.5};
\addplot[curveblue, thick, domain=0:3.2] {-(0.45*x+0.5)};
\draw[black, fill=areafill!40] (axis cs:1.5,0) ellipse [x radius=3pt, y radius=24pt];
\draw[black, fill=areafill!60] (axis cs:1.75,0) ellipse [x radius=3pt, y radius=26pt];
\draw[tangentred, very thick] (axis cs:1.75,0) -- (axis cs:1.75,1.288) node[midway, right, font=\footnotesize] {$f(x)$};
\fill[black] (axis cs:1.75,0) circle (1.5pt);
\draw[tangentred, line width=1.8pt] (axis cs:1.5,1.175) -- (axis cs:1.75,1.288)
  node[midway, above left, font=\footnotesize] {$\dd s$};
\draw[->] (axis cs:1.625,-1.8) node[below, font=\footnotesize] {band of slant width $\dd s$} -- (axis cs:1.625,-1.35);
\end{axis}
\end{tikzpicture}
\caption{Revolving $y=f(x)$ about the $x$-axis. A short arc of length $\dd s$ sweeps
out a band of radius $f(x)$; its area is $2\pi f(x)\,\dd s$.}
\label{ch26:fig-surface}
\end{figure}

\begin{whybox}
\textbf{Why $2\pi f(x)\,\dd s$ and not $2\pi f(x)\,\dd x$?} A natural instinct is
to multiply the circumference $2\pi f(x)$ by the horizontal width $\dd x$. That is
wrong, and the reason is geometric. The little band is not a vertical cylinder; it
is a slanted \emph{frustum} of a cone -- the shape of a lampshade or a drinking cup
with the bottom cut open -- because the curve is rising. Its slant width is the
\emph{arc length} $\dd s$, which is longer than $\dd x$ whenever the curve has any
slope. Using $\dd x$ would systematically undercount the surface -- it ignores the
tilt. The correct area of a thin frustum band of radius $r$ and slant width
$\ell$ is $2\pi r\,\ell$; here $r=f(x)$ and $\ell = \dd s$. So the $\dd s$ is not a
decoration; it is the whole point. (This is exactly the surface-area analogue of
why arc length needs the radical.)
\end{whybox}

\begin{example}[Surface area of a sphere]
\label{ch26:ex-sphere-surface}
Revolve $y = \sqrt{r^2-x^2}$, from $x=-r$ to $x=r$, about the $x$-axis. We have
$f(x)=\sqrt{r^2-x^2}$ and
\[
  f'(x) = \frac{-x}{\sqrt{r^2-x^2}}, \qquad
  1+[f'(x)]^2 = 1 + \frac{x^2}{r^2-x^2} = \frac{r^2}{r^2-x^2}.
\]
The radical is therefore $\sqrt{1+[f']^2} = \dfrac{r}{\sqrt{r^2-x^2}}$, giving
\[
  \dd S = 2\pi\sqrt{r^2-x^2}\cdot\frac{r}{\sqrt{r^2-x^2}}\,\dd x = 2\pi r\,\dd x.
\]
The square roots cancel cleanly, leaving a constant. Then
\[
  S = \int_{-r}^{r} 2\pi r\,\dd x = 2\pi r\cdot(2r) = 4\pi r^2.
\]
The surface area of a sphere is $4\pi r^2$.
\end{example}

\begin{aestheticbox}
The cancellation in that calculation reflects an exact geometric theorem known
to Archimedes. Because $\dd S = 2\pi r\,\dd x$ is constant with respect to $x$, slices of a
sphere between parallel planes have surface area proportional strictly to the axial
thickness $\Delta x$, regardless of where along the axis the slice is taken. A spherical
cap near a pole and a band near the equator have equal surface areas provided they
span the same axial height. This is \emph{Archimedes' hat-box theorem}: a sphere and its
circumscribed cylinder, cut by two planes perpendicular to the cylinder's axis, have equal
lateral surface areas between those planes. This property provides the foundation for the
Lambert cylindrical equal-area projection in cartography.
\end{aestheticbox}

\section{Work: Accumulating Force Over Distance}
\label{ch26:work}

We leave geometry and integrate a physical process. In physics, \emph{work} is
force times distance -- but only when the force is constant. What if the force
changes as the object moves? Slice the motion.

\textbf{Slice.} Break the path into tiny displacements $\dd x$.

\textbf{Measure.} Over a displacement so small that the force $F(x)$ barely
changes, the work done is (force)$\times$(distance) $= F(x)\,\dd x$.

\textbf{Sum.}
\[
  \boxed{\,W = \int_a^b F(x)\,\dd x.\,}
\]
Work is the accumulation of force over distance. This is the prototype of every
``integrate a varying rate'' problem in science.

\begin{example}[Stretching a spring -- Hooke's law]
\label{ch26:ex-spring}
A spring resists stretching with a force proportional to how far it is stretched:
$F(x) = kx$ (Hooke's law), where $k$ is the spring constant. How much work does it
take to stretch the spring from its natural length to a displacement $x=d$?

The force is not constant -- it grows as you pull -- so you must integrate:
\[
  W = \int_0^d kx\,\dd x = k\cdot\frac{x^2}{2}\Big|_0^d = \frac{1}{2}kd^2.
\]
The $\tfrac12$ appears for the same reason it always does: you are summing a
linearly growing quantity, and the average of a quantity that ramps from $0$ to
$kd$ is $\tfrac12 kd$. Work = average force times distance = $\tfrac12 kd\cdot d$.
The integral and the averaging picture agree.
\end{example}

\begin{connectionbox}
\textbf{Work across an infinite horizon.} What if the force extends outward indefinitely?
In Chapter~25 you learned how to evaluate improper integrals over unbounded intervals.
Consider Newton's law of universal gravitation: the attractive force on a spacecraft of mass
$m$ at distance $r$ from the center of a planet of mass $M$ is $F(r) = \frac{G M m}{r^2}$.
How much work is required to haul the spacecraft from the planet's surface at radius $R$
completely free of its gravitational field?

The displacement stretches to infinity, so the work is an improper integral:
\[
  W = \int_R^{\infty} \frac{G M m}{r^2}\,\dd r
    = \lim_{b\to\infty}\left[ -\frac{G M m}{r} \right]_R^b
    = 0 - \left(-\frac{G M m}{R}\right)
    = \frac{G M m}{R}.
\]
Because the $1/r^2$ force drops off quickly enough, the infinite stretch of space costs
only a finite expenditure of energy. Setting this work equal to the spacecraft's initial
kinetic energy $\tfrac{1}{2}m v_{\text{esc}}^2$ yields the famous escape velocity formula
$v_{\text{esc}} = \sqrt{2GM/R}$. Because gravity falls off as $1/r^2$ ($p=2 > 1$),
the improper integral converges and a finite launch speed suffices to escape.
\end{connectionbox}

\subsection{Pumping problems: when the slices are layers of fluid}

The richest work problems involve pumping liquid out of a tank. Here the slicing is
subtle and worth dwelling on, because \emph{each slice is a different amount of
stuff that has to move a different distance}.

The setup: a tank holds water; you must pump it all up and out over the top rim.
Slice the water into thin horizontal layers. Consider one layer at height $y$,
thickness $\dd y$.

\textbf{How much does this layer weigh?} Its volume is (cross-sectional area at
height $y$) times thickness, $A(y)\,\dd y$. Multiply by the weight density of the
fluid $\rho g$ (for water, about $9800\ \text{N/m}^3$, or in US units the weight
density is about $62.4\ \text{lb/ft}^3$). So the layer's weight is
$\rho g\, A(y)\,\dd y$.

\textbf{How far must this layer move?} If the water exits at height $H$, this layer
must rise a distance $(H - y)$.

\textbf{Work to lift this one layer} = weight $\times$ distance:
\[
  \dd W = \rho g\, A(y)\,(H-y)\,\dd y.
\]

\textbf{Sum over all layers:}
\[
  \boxed{\,W = \rho g\int_{y_{\min}}^{y_{\max}} A(y)\,(H-y)\,\dd y.\,}
\]

\begin{figure}[H]
\centering
\begin{tikzpicture}
\draw[thick] (-1.6,3)--(-0.8,0)--(0.8,0)--(1.6,3);
\draw[thick, dashed] (-1.6,3)--(1.6,3);
\fill[areafill!55] (-1.18,1.3) rectangle (1.18,1.5);
\draw (-1.18,1.3)--(1.18,1.3); \draw(-1.18,1.5)--(1.18,1.5);
\node[right] at (1.7,1.4) {\footnotesize layer at height $y$, thickness $\dd y$};
\draw[->, gray] (0,1.4)--(0,3.2);
\node[left, gray] at (0,2.4) {\footnotesize lift $H-y$};
\node at (0,-0.35) {\footnotesize $y=0$};
\node at (0,3.3) {\footnotesize $y=H$ (rim)};
\end{tikzpicture}
\caption{Pumping a tank. A horizontal layer at height $y$ weighs $\rho g A(y)\,\dd
y$ and must be lifted a distance $H-y$ to clear the rim.}
\label{ch26:fig-pump}
\end{figure}

\begin{trapbox}
\textbf{The two quantities that vary are easy to confuse.} In a pumping problem
\emph{two} things depend on the slice: how much water is in the layer (through
$A(y)$) and how far it travels (through $H-y$). Beginners get one and forget the
other -- using a constant cross-section for a cone, or forgetting that lower layers
must travel farther. Always write both factors explicitly and ask, for a test
layer, ``is this the weight? is this the distance?'' Also: the lift distance is
measured to the \emph{exit point}, which may be above the tank's rim if there is a
spout. Read the geometry carefully.
\end{trapbox}

\begin{example}[Pumping a cylindrical tank]
\label{ch26:ex-pump-cyl}
A cylindrical tank, radius $2\ \text{m}$ and height $5\ \text{m}$, is full of water.
Find the work to pump all the water out over the top rim. Use $\rho g = 9800\
\text{N/m}^3$.

Put $y=0$ at the bottom, $y=5$ at the rim, so $H=5$. The cross-section is a constant
disk of area $A(y) = \pi(2)^2 = 4\pi$. A layer at height $y$ must rise $(5-y)$. So
\[
  W = 9800\int_0^5 4\pi\,(5-y)\,\dd y
    = 9800\cdot 4\pi\left[5y - \frac{y^2}{2}\right]_0^5.
\]
The bracket is $25 - \tfrac{25}{2} = \tfrac{25}{2}$. Hence
\[
  W = 9800\cdot 4\pi\cdot\frac{25}{2} = 9800\cdot 50\pi = 490000\,\pi
    \approx 1.54\times 10^{6}\ \text{J}.
\]
About $1.5$ million joules. Notice that even though every layer weighs the same
(constant cross-section), they do \emph{not} all do the same work, because the
lower layers must be lifted farther -- the $(5-y)$ factor does the bookkeeping.
\end{example}

\begin{labbox}
Redo Example~\ref{ch26:ex-pump-cyl} but pump the water to a spout $1\ \text{m}$
above the rim, so $H=6$. You should only need to change one number and re-evaluate
$\int_0^5 (6-y)\,\dd y = [6y - y^2/2]_0^5 = 30 - 12.5 = 17.5$. Confirm
$W = 9800\cdot 4\pi\cdot 17.5 = 686000\,\pi \approx 2.16\times 10^6\ \text{J}$.
Notice how much extra work that one meter of extra lift costs -- and why: every
single layer pays the surcharge.
\end{labbox}

\section{Average Value of a Function}
\label{ch26:average}

How do you average infinitely many numbers -- all the values $f(x)$ takes as $x$
ranges over $[a,b]$? The average of finitely many numbers is their sum divided by
their count. For a continuous spread of values, the sum becomes an integral and the
count becomes the length of the interval:
\[
  \boxed{\,\bar f = \frac{1}{b-a}\int_a^b f(x)\,\dd x.\,}
\]

Where does this come from? Sample $f$ at $n$ evenly spaced points and average:
\[
  \frac{1}{n}\sum_{k=1}^n f(x_k)
  = \frac{1}{b-a}\sum_{k=1}^n f(x_k)\,\frac{b-a}{n}
  = \frac{1}{b-a}\sum_{k=1}^n f(x_k)\,\Delta x.
\]
We slipped $\Delta x = (b-a)/n$ in and out. As $n\to\infty$ the sum becomes the
integral, giving the boxed formula. So the average value is the integral spread
evenly over the interval.

\begin{intuitionbox}
The average value $\bar f$ is the height of the rectangle, sitting on $[a,b]$, that
has the \emph{same area} as the region under $f$. You are flattening the bumpy
graph into a level one without changing the total area underneath. Rearranged,
$\bar f\,(b-a) = \int_a^b f$: average times width equals area. That is the cleanest
way to remember it.
\end{intuitionbox}

\begin{figure}[H]
\centering
\begin{tikzpicture}
\begin{axis}[
  width=9cm, height=5cm,
  xmin=0, xmax=4, ymin=0, ymax=3.2,
  axis lines=left, xlabel={$x$}, ylabel={$y$},
  xtick={0.5,3.5}, xticklabels={$a$,$b$}, ytick={1.6}, yticklabels={$\bar f$},
  samples=100,
]
\addplot[name path=f, curveblue, very thick, domain=0.5:3.5]
  {1.6 + 1.1*sin(deg(2*x))};
\addplot[name path=axisline, draw=none, domain=0.5:3.5] {0};
\addplot[areafill!30] fill between[of=f and axisline];
\addplot[tangentred, thick, dashed, domain=0.5:3.5] {1.6};
\end{axis}
\end{tikzpicture}
\caption{The average value $\bar f$ is the height of the level rectangle over
$[a,b]$ whose area equals the area under $f$. The wiggles above the dashed line
exactly fill the gaps below it.}
\label{ch26:fig-average}
\end{figure}

\begin{example}[Average value]
\label{ch26:ex-average}
Find the average value of $f(x) = x^2$ on $[0,3]$.
\[
  \bar f = \frac{1}{3-0}\int_0^3 x^2\,\dd x
         = \frac13\cdot\frac{x^3}{3}\Big|_0^3
         = \frac13\cdot 9 = 3.
\]
The function climbs from $0$ to $9$, but because it spends most of the interval at
smaller values, its average is only $3$, well below $9/2$, the naive average of the
two endpoint values $0$ and $9$. Averages of nonlinear functions are not the average
of the endpoints -- the shape matters.
\end{example}

\begin{connectionbox}
The \emph{Mean Value Theorem for Integrals} says that a continuous $f$ actually
\emph{attains} its average value: there is some $c\in[a,b]$ with $f(c)=\bar f$.
Geometrically, the curve must cross the level $\bar f$ line somewhere -- it cannot
be entirely above its own average. This is the integral cousin of the Mean Value
Theorem for derivatives you met earlier; both say a continuous quantity must,
somewhere, equal its average behavior.
\end{connectionbox}

\section{A Preview of Probability Density}
\label{ch26:probability}

We finish with the application that integration was, in a sense, reinvented for in
the twentieth century: \emph{continuous probability}. The slicing idea here is so
natural that probability almost forces it on you.

How do you describe the probability of a quantity that can take any real value --
a person's exact height, the precise time a bus arrives, the energy of a molecule?
You cannot assign positive probability to each single value, because there are
uncountably many and they would sum to more than $1$. The probability of the height
being \emph{exactly} $1.8000\ldots$ meters is zero. So instead you describe
probability as something \emph{spread out}, a density, and you accumulate it by
integrating.

A \textbf{probability density function} (pdf) is a function $p(x) \ge 0$ with total
integral $1$:
\[
  \int_{-\infty}^{\infty} p(x)\,\dd x = 1.
\]
Because the real line is unbounded, this total probability is itself an improper
integral of the exact kind you explored in Chapter~25: the probability of all possibilities
combined must converge to certainty.

The meaning of $p$ is exactly the slicing meaning: the probability that the quantity
lands in a tiny window $[x, x+\dd x]$ is
\[
  \dd(\text{prob}) = p(x)\,\dd x,
\]
and the probability of landing in any range $[a,b]$ is the accumulation of those
slices (Figure~\ref{ch26:fig-prob-density}):
\[
  \boxed{\,P(a \le X \le b) = \int_a^b p(x)\,\dd x.\,}
\]
Probability is area under the density curve. That is the whole conceptual content.

\begin{figure}[H]
\centering
\begin{tikzpicture}
\begin{axis}[
  width=10cm, height=5.2cm,
  xmin=-0.5, xmax=4.5, ymin=0, ymax=0.6,
  axis lines=left, xlabel={$x$}, ylabel={$p(x)$},
  xtick={1.2, 2.8}, xticklabels={$a$, $b$},
  ytick=\empty, samples=100,
  clip=false,
]
\addplot[name path=curve, curveblue, very thick, domain=0:4.2]
  {0.52 * exp(-0.5*((x-2)/0.9)^2)};
\addplot[name path=axisline, draw=none, domain=0:4.2] {0};
\addplot[areafill!45] fill between[of=curve and axisline, soft clip={domain=1.2:2.8}];
\node at (axis cs:1.7, 0.22) {\small $P(a \le X \le b)$};
\fill[quotientorange!60] (axis cs:2.28, 0) rectangle (axis cs:2.42, 0.47);
\draw[black, thick] (axis cs:2.28, 0) rectangle (axis cs:2.42, 0.47);
\node[above] at (axis cs:2.35, 0.48) {\footnotesize $p(x)$};
\draw[<->] (axis cs:2.28, -0.04) -- (axis cs:2.42, -0.04)
  node[midway, below, font=\tiny] {$\dd x$};
\node[curveblue, anchor=south west] at (axis cs:2.8, 0.48) {\small $\dd P = p(x)\,\dd x$};
\draw[->, curveblue] (axis cs:2.8, 0.49) -- (axis cs:2.44, 0.49);
\end{axis}
\end{tikzpicture}
\caption{A continuous probability density $p(x)$. The probability of landing in the
window $[a,b]$ is the accumulated area under the curve. A single infinitesimal slice
carries probability $\dd P = p(x)\,\dd x$.}
\label{ch26:fig-prob-density}
\end{figure}

\begin{happeningbox}
A density is not a probability -- it is a probability \emph{per unit length}, just
as mass density is mass per unit volume. By itself $p(x)$ can even exceed $1$; what
must stay $\le 1$ is the \emph{area} $p(x)\,\dd x$ accumulated over a region. This
is why the units work and why ``the probability of exactly $x$'' is zero: a single
point has width $\dd x = 0$, so it carries area zero. Probability lives in the
\emph{accumulation}, never in the point. To get a real probability you must
integrate the density across a window of nonzero width -- you must collect a band
of slices.
\end{happeningbox}

\begin{example}[The uniform and exponential densities]
\label{ch26:ex-density}
\textbf{Uniform on $[0,L]$.} A bus is equally likely to arrive at any time in
$[0,L]$. The density is constant, $p(x) = 1/L$ on $[0,L]$ and $0$ outside. Check:
$\int_0^L \tfrac1L\,\dd x = 1$. The probability the bus comes in the first third is
$\int_0^{L/3}\tfrac1L\,\dd x = \tfrac13$ -- area of a rectangle, exactly as
intuition demands.

\textbf{Exponential.} The waiting time for a random event (a radioactive decay, a
phone call) often follows $p(x) = \lambda e^{-\lambda x}$ for $x\ge 0$. First check
it is a valid density. Because the domain $[0,\infty)$ is unbounded, this is an
improper integral of the first kind (Chapter~25):
\[
  \int_0^{\infty}\lambda e^{-\lambda x}\,\dd x
  = \lim_{b\to\infty}\Bigl[-e^{-\lambda x}\Bigr]_0^{b}
  = \lim_{b\to\infty}\bigl(-e^{-\lambda b} - (-1)\bigr)
  = 0 + 1 = 1.
\]
The probability the event happens before time $t$ is
\[
  P(X \le t) = \int_0^t \lambda e^{-\lambda x}\,\dd x
             = 1 - e^{-\lambda t}.
\]
At $t=0$ this is $0$; as $t\to\infty$ it climbs to $1$. The accumulation of density
sweeps from impossible to certain.
\end{example}

\begin{applicationbox}
The \textbf{expected value} (mean) of a continuous random variable is, once again,
a slicing problem. To average $x$ over the density, weight each value $x$ by the
probability $p(x)\,\dd x$ sitting on it, and sum:
\[
  \mathbb{E}[X] = \int_{-\infty}^{\infty} x\,p(x)\,\dd x.
\]
This is the continuous analogue of $\sum (\text{value})\times(\text{probability})$,
and it is also exactly the formula for the \emph{center of mass} of a rod whose
mass density is $p$ -- the balance point. Probability theory and mechanics share
the integral because they share the slice: in both, you accumulate a quantity
weighted by how much ``stuff'' (probability, or mass) sits at each location.

For the exponential density on $[0,\infty)$, we evaluate this improper integral
using integration by parts (Chapter~23) with $u = x$ and $\dd v = \lambda e^{-\lambda x}\,\dd x$
(so $\dd u = \dd x$ and $v = -e^{-\lambda x}$):
\[
  \mathbb{E}[X] = \int_0^{\infty} x\cdot \lambda e^{-\lambda x}\,\dd x
  = \lim_{b\to\infty}\left( \Bigl[-x e^{-\lambda x}\Bigr]_0^b + \int_0^b e^{-\lambda x}\,\dd x \right)
  = 0 + \lim_{b\to\infty}\left[ -\frac{1}{\lambda}e^{-\lambda x} \right]_0^b
  = \frac{1}{\lambda}.
\]
The famous mean waiting time $1/\lambda$ unites improper integrals, integration
by parts, and physical accumulation in a single calculation.
\end{applicationbox}

\section{The Common Thread}
\label{ch26:thread}

Look back at everything we built. We solved eight very different problems --
area, volume several ways, arc length, surface area, work, pumping, average value,
probability -- and we never once invented a genuinely new idea. Every single
result came from the same three-step move:

\begin{center}
\begin{tabular}{@{}lll@{}}
\toprule
\textbf{Quantity} & \textbf{One slice $\dd Q$} & \textbf{Accumulated} \\
\midrule
Area between curves & $(f-g)\,\dd x$ & $\int_a^b (f-g)\,\dd x$ \\
Volume (slicing) & $A(x)\,\dd x$ & $\int_a^b A(x)\,\dd x$ \\
Volume (disks) & $\pi[f]^2\,\dd x$ & $\pi\int_a^b [f]^2\,\dd x$ \\
Volume (washers) & $\pi([f]^2-[g]^2)\,\dd x$ & $\pi\int_a^b([f]^2-[g]^2)\,\dd x$ \\
Volume (shells) & $2\pi x\,f(x)\,\dd x$ & $\int_a^b 2\pi x f(x)\,\dd x$ \\
Arc length & $\sqrt{1+[f']^2}\,\dd x$ & $\int_a^b\sqrt{1+[f']^2}\,\dd x$ \\
Surface area & $2\pi f\sqrt{1+[f']^2}\,\dd x$ & $\int_a^b 2\pi f\sqrt{1+[f']^2}\,\dd x$ \\
Work & $F(x)\,\dd x$ & $\int_a^b F(x)\,\dd x$ \\
Pumping & $\rho g A(y)(H-y)\,\dd y$ & $\rho g\int A(y)(H-y)\,\dd y$ \\
Average value & --- & $\frac{1}{b-a}\int_a^b f\,\dd x$ \\
Probability & $p(x)\,\dd x$ & $\int_a^b p(x)\,\dd x$ \\
\bottomrule
\end{tabular}
\end{center}

The middle column is the only one that ever required thought, and it always came
straight from a picture: a rectangle, a coin, a shell, a triangle, a layer of
water. \emph{Find the slice, and the integral writes itself.} That is the entire
craft of applied integration, and it will carry you through multiple integrals,
center of mass, fluid pressure, and far beyond. Whenever you meet a new quantity to
compute, do not hunt for a formula. Ask the one question this chapter is built on:
\emph{what am I accumulating, and what does one slice contribute?}

\begin{aestheticbox}
In classical geometry, Archimedes and Eudoxus had to invent a separate exhaustion argument for
each shape -- one construction for a parabolic segment, another for a sphere, another for a spiral.
Writing $\dd Q$ for a single thin slice and integrating replaces all of those separate geometric
arguments with one uniform calculation.
\end{aestheticbox}

\section{Limitations and the road ahead}

\begin{connectionbox}
The slicing method expresses many geometric and physical quantities as integrals:
areas, volumes of revolution, arc lengths, surface areas, work, and probability densities.

However, setting up an integral does not guarantee that its antiderivative is elementary.
Many natural geometry and physics problems lead to integrals that cannot be evaluated
in closed form:
\begin{itemize}
  \item \emph{Arc length of an ellipse:} Slicing the ellipse $\frac{x^2}{a^2} + \frac{y^2}{b^2} = 1$
  yields the perimeter
  \[
    L = 4a \int_0^{\pi/2} \sqrt{1 - e^2 \sin^2\theta}\,\dd\theta,
  \]
  where $e = \sqrt{1 - b^2/a^2}$ is the eccentricity. This is a \emph{complete elliptic
  integral of the second kind}, which has no elementary antiderivative.
  \item \emph{Pendulum period:} Computing the exact period of a pendulum swinging from an
  initial angle $\theta_0$ yields a \emph{complete elliptic integral of the first kind}:
  \[
    T = 4\sqrt{\frac{L_p}{g}} \int_0^{\pi/2} \frac{\dd\phi}{\sqrt{1 - \sin^2(\theta_0/2)\sin^2\phi}}.
  \]
\end{itemize}

When closed-form formulas are unavailable, we evaluate integrals using \emph{numerical integration}
(Chapter~27) or approximate them with series expansions (Chapter~30). In Chapter~27, we
develop algorithms that approximate definite integrals directly from sample values to high accuracy.
\end{connectionbox}

\section{Exercises}
\label{ch26:exercises}

\begin{exercise}[Mental check: round-trip work]
A worker lifts a crate of mass $m$ vertically from the floor to a shelf at height $h$,
and later lowers it back to the floor. Gravity exerts a constant downward force of magnitude
$mg$. (a) How much work is done by gravity as the crate is lifted? (b) How much work is done
by gravity as the crate is lowered? (c) What is the total work done by gravity over the round
trip? Explain why the signs must differ.
\end{exercise}

\begin{exercise}[Mental check: disk vs.\ shell selection]
Consider the region bounded by $y = e^{-x^2}$, $y = 0$, $x = 0$, and $x = 1$, revolved
about the $y$-axis. Would you choose the disk/washer method or the shell method to set up
the volume? Explain your choice in one sentence without computing either integral.
(Hint: consider the antiderivatives of $x e^{-x^2}$ versus $\sqrt{-\ln y}$.)
\end{exercise}

\begin{exercise}[Area, both directions]
Find the area of the region bounded by $y = \sqrt{x}$ and $y = x^2$. First by
vertical slices, then by horizontal slices, and confirm the answers agree.
\end{exercise}

\begin{exercise}[Crossing curves]
Find the total area enclosed between $y=\sin x$ and $y=\cos x$ from $x=0$ to
$x=\pi/2$. (The curves cross at $x=\pi/4$; remember to use top-minus-bottom on each
piece, so the answer is positive.)
\end{exercise}

\begin{exercise}[Volume by slicing, non-circular cross sections]
A solid has as its base the region between $y=0$ and $y = 1-x^2$ for $-1\le x\le 1$.
Every cross-section perpendicular to the $x$-axis is a square sitting on the base.
Set up and evaluate $\int A(x)\,\dd x$ to find the volume. (Here the side of the
square at position $x$ is $1-x^2$, so $A(x) = (1-x^2)^2$.)
\end{exercise}

\begin{exercise}[Disk vs.\ shell]
Let $R$ be the region under $y = x^2$ from $x=0$ to $x=1$. Find the volume when $R$
is revolved (a) about the $x$-axis, and (b) about the $y$-axis. For each, decide
whether disks/washers or shells gives the easier integral, and explain your choice
in one sentence.
\end{exercise}

\begin{exercise}[Arc length]
Find the arc length of $y = \tfrac{1}{3}\bigl(x^2+2\bigr)^{3/2}$ from $x=0$ to
$x=2$. (Compute $f'$, simplify $1+[f']^2$ to a perfect square, and watch the
radical disappear.)
\end{exercise}

\begin{exercise}[Surface area]
The curve $y=\sqrt{x}$ from $x=0$ to $x=1$ is revolved about the $x$-axis. Set up
the surface-area integral $\int 2\pi f\sqrt{1+[f']^2}\,\dd x$, simplify the
integrand, and evaluate it.
\end{exercise}

\begin{exercise}[Work, varying force]
A chain $10\ \text{m}$ long weighs $5\ \text{N/m}$ and hangs from the top of a
building. How much work is required to haul the entire chain to the top? (Hint:
slice the chain by length; a piece at depth $x$ weighs $5\,\dd x$ and must rise a
distance $x$, so $\dd W = 5x\,\dd x$.)
\end{exercise}

\begin{exercise}[Pumping a cone]
A tank in the shape of an inverted cone (point down) has height $4\ \text{m}$ and
top radius $2\ \text{m}$, and is full of water. Set up and evaluate the work to pump
all the water to the top rim. (By similar triangles the radius at height $y$ is
$y/2$, so $A(y) = \pi(y/2)^2$; the lift is $4-y$. Use $\rho g = 9800\
\text{N/m}^3$.)
\end{exercise}

\begin{exercise}[Average value, and where it is attained]
Find the average value of $f(x) = \sin x$ on $[0,\pi]$. Then find the value(s) of
$c\in[0,\pi]$ where $f(c)$ equals that average, illustrating the Mean Value Theorem
for Integrals.
\end{exercise}

\begin{exercise}[Probability density]
A continuous random variable has density $p(x) = c\,x(1-x)$ on $[0,1]$ and zero
elsewhere. (a) Find $c$ so that $\int_0^1 p = 1$. (b) Compute $P(X \le 1/2)$.
(c) Compute the mean $\mathbb{E}[X] = \int_0^1 x\,p(x)\,\dd x$, and explain why
symmetry could have told you the answer in advance.
\end{exercise}

\begin{exercise}[Conceptual: the slice is everything]
In one short paragraph each, state ``what one slice contributes'' ($\dd Q$) for:
the area between two curves, a washer, a cylindrical shell, and the work to pump a
fluid. You should be able to do this from the pictures alone, without looking up a
formula -- that is the point of the chapter.
\end{exercise}

\begin{exercise}[Proof: shells and washers must agree]
Consider the region under $y=f(x)$ from $x=0$ to $x=b$ (with $f$ decreasing,
$f(b)=0$) revolved about the $y$-axis. Write the shell-method volume
$V = \int_0^b 2\pi x f(x)\,\dd x$. Using integration by parts with $u=f(x)$ and
$\dd v = 2\pi x\,\dd x$, transform this into an integral over $y$ and identify the
result as the washer/disk volume. Conclude that the two methods are not two facts
but one, related by integration by parts. (This is a satisfying capstone: the
algebra of parts is exactly the geometry of re-slicing.)
\end{exercise}
